% !TEX root = ../../ctfp-print.tex

\lettrine[lhang=0.17]{圏}{論の本を締めくくるのに} 良い場所などというものはない。学ぶべきことは常にある。圏論は広大な分野だ。同時に、同じテーマ、概念、パターンが何度も繰り返し現れることは明らかだ。すべての概念はカン拡張であるという格言があり、実際にカン拡張を使って極限、余極限、随伴、モナド、米田の補題などを導出できる。圏そのものの概念はあらゆる抽象化のレベルで現れ、モノイドとモナドの概念も同様だ。どれが最も基本的なのだろうか？実はそれらはすべて相互に関連しており、一方から他方へと抽象化の終わりなきサイクルを形成している。これらの相互接続を示すことが、この本を締めくくる良い方法だと判断した。

\section{双圏}

圏論で最も難しい側面の一つは、視点を絶えず切り替えることだ。たとえば集合の圏を考えてみよう。私たちは要素によって集合を定義することに慣れている。空集合には要素がない。単元集合には一つの要素がある。二つの集合のデカルト積はペアの集合である、などなど。しかし圏 $\Set$ について語るとき、私は集合の中身を忘れ、代わりに集合間の射（矢）に集中するよう求めた。$\Set$ における特定の普遍的構成が要素によって何を意味するかを見るために、時折カバーの下を覗くことは許されていた。終対象は一つの要素を持つ集合であることが分かった、などだ。しかしこれらは単なる健全性確認に過ぎなかった。

関手は圏の写像として定義される。写像を圏における射と考えるのは自然だ。関手は圏の圏（大きさに関する問題を避けるなら、小圏の）における射であることが分かった。関手を矢として扱うことで、圏の内部（その対象と射）への作用に関する情報を失う。ちょうど関数を $\Set$ における矢として扱うとき、集合の要素への作用に関する情報を失うのと同様だ。しかし任意の二圏間の関手もまた圏を形成する。今度は、ある圏において矢であったものを別の圏で対象として考えることが求められる。関手圏では、関手が対象であり自然変換が射だ。同じものがある圏では矢であり、別の圏では対象であり得ることが分かった。対象を名詞とし矢を動詞とする素朴な見方は成り立たない。

二つの視点を切り替える代わりに、それらを一つに統合しようとすることができる。これが $\cat{2}$-圏の概念を得る方法で、対象を $0$-セル、射を $1$-セル、射の間の射を $2$-セルと呼ぶ。

\begin{figure}[H]
  \centering
  \includegraphics[width=0.35\textwidth]{images/twocat.png}
  \caption{$0$-セル $a, b$；$1$-セル $f, g$；および $2$-セル $\alpha$。}
\end{figure}

\noindent
圏の圏 $\Cat$ はすぐに思い浮かぶ例だ。$0$-セルとして圏、$1$-セルとして関手、$2$-セルとして自然変換を持つ。$\cat{2}$-圏の法則は、任意の二つの $0$-セル間の $1$-セルが圏を形成することを教えてくれる（言い換えれば、$\cat{C}(a, b)$ はhom集合ではなくhom圏である）。これは、任意の二圏間の関手が関手圏を形成するという以前の主張とうまく合致する。

特に、任意の $0$-セルからそれ自身への $1$-セルもまた圏、hom圏 $\cat{C}(a, a)$ を形成するが、その圏にはさらに多くの構造がある。$\cat{C}(a, a)$ のメンバーは $\cat{C}$ における矢として、あるいは $\cat{C}(a, a)$ における対象として見ることができる。矢として、それらは互いに合成できる。しかし対象として見るとき、合成は対象のペアから対象への写像になる。実際、これはテンソル積に非常によく似ている。このテンソル積には単位元がある：恒等 $1$-セルだ。任意の $\cat{2}$-圏において、hom圏 $\cat{C}(a, a)$ は $1$-セルの合成によって定義されるテンソル積を持つモノイダル圏に自動的になることが分かる。結合律と単位律は対応する圏の法則から自然に導出される。

$\cat{2}$-圏 $\Cat$ の標準的な例でこれが何を意味するか見てみよう。hom圏 $\Cat(a, a)$ は $a$ 上の自己関手の圏だ。自己関手の合成がテンソル積の役割を果たす。恒等関手がこの積に関する単位元だ。以前に自己関手がモノイダル圏を形成することを見た（モナドの定義でこの事実を使った）が、今ではこれがより一般的な現象であることが分かる：任意の $\cat{2}$-圏における自己 $1$-セルはモノイダル圏を形成する。後でモナドを一般化するときにこれに戻ってくる。

一般的なモノイダル圏では、モノイドの法則が厳密に満たされることを主張しなかったことを思い出すかもしれない。単位律と結合律が同型まで満たされれば十分なことが多かった。$\cat{2}$-圏では、$\cat{C}(a, a)$ におけるモノイダル法則は $1$-セルの合成法則から導かれる。これらの法則は厳密なので、常に厳密モノイダル圏を得ることになる。しかしこれらの法則を緩めることも可能だ。たとえば、恒等 $1$-セル $\idarrow[a]$ と別の $1$-セル $f \Colon a \to b$ との合成が、$f$ と等しいのではなく同型であると言うことができる。$1$-セルの同型は $2$-セルを使って定義される。言い換えれば、逆元を持つ $2$-セル：
\[\rho \Colon f \circ \idarrow[a] \to f\]
が存在する。

\begin{figure}[H]
  \centering
  \includegraphics[width=0.35\textwidth]{images/bicat.png}
  \caption{双圏における単位律は同型（可逆な $2$-セル $\rho$）まで成り立つ。}
\end{figure}

\noindent
左単位律と結合律についても同様にできる。このような緩められた $\cat{2}$-圏は双圏と呼ばれる（ここでは省略するが、追加のコヒーレンス条件がある）。

予想通り、双圏における自己 $1$-セルは非厳密な法則を持つ一般的なモノイダル圏を形成する。

双圏の興味深い例はスパンの圏だ。二つの対象 $a$ と $b$ の間のスパンとは、対象 $x$ と射のペア：
\begin{gather*}
  f \Colon x \to a \\
  g \Colon x \to b
\end{gather*}
である。

\begin{figure}[H]
  \centering
  \includegraphics[width=0.35\textwidth]{images/span.png}
\end{figure}

\noindent
スパンを圏積の定義で使ったことを思い出すかもしれない。ここでは、スパンを双圏における $1$-セルとして見たい。最初のステップはスパンの合成を定義することだ。次のように接続するスパンがあるとしよう：
\begin{gather*}
  f' \Colon y \to b \\
  g' \Colon y \to c
\end{gather*}

\begin{figure}[H]
  \centering
  \includegraphics[width=0.5\textwidth]{images/compspan.png}
\end{figure}

\noindent
合成は、頂点 $z$ を持つ第三のスパンになる。最も自然な選択は、$f'$ に沿った $g$ の引き戻しだ。引き戻しとは、二つの射を持つ対象 $z$ であることを思い出そう：
\begin{align*}
  h  & \Colon z \to x \\
  h' & \Colon z \to y
\end{align*}
が次を満たす：
\[g \circ h = f' \circ h'\]
これはそのようなすべての対象の中で普遍的なものだ。

\begin{figure}[H]
  \centering
  \includegraphics[width=0.5\textwidth]{images/pullspan.png}
\end{figure}

\noindent
今は集合の圏上のスパンに集中しよう。その場合、引き戻しはデカルト積 $x \times y$ からのペアの集合 $(p, q)$ で次を満たすものだ：
\[g\ p = f'\ q\]
同じ端点を共有する二つのスパン間の射は、適切な三角形が可換になるような頂点間の射 $h$ として定義される。

\begin{figure}[H]
  \centering
  \includegraphics[width=0.4\textwidth]{images/morphspan.png}
  \caption{$\cat{Span}$ における $2$-セル。}
\end{figure}

\noindent
要約すると、双圏 $\cat{Span}$ において：$0$-セルは集合、$1$-セルはスパン、$2$-セルはスパンの射だ。恒等 $1$-セルは、三つの対象がすべて同じで二つの射が恒等射である退化したスパンだ。

双圏のもう一つの例を以前に見た：\hyperref[ends-and-coends]{プロ関手}の双圏 $\cat{Prof}$ で、$0$-セルは圏、$1$-セルはプロ関手、$2$-セルは自然変換だ。プロ関手の合成は余端によって与えられた。

\section{モナド}

今では、モナドを自己関手の圏におけるモノイドとして定義することにかなり慣れているはずだ。自己関手の圏が双圏 $\Cat$ における自己 $1$-セルのhom圏の一つに過ぎないという新しい理解でこの定義を再訪しよう。それがモノイダル圏であることは分かっている：テンソル積は自己関手の合成から来る。モノイドはモノイダル圏における対象として定義される——ここでは自己関手 $T$ になる——二つの射と共に。自己関手間の射は自然変換だ。一方の射はモノイダル単位——恒等自己関手——を $T$ に写す：
\[\eta \Colon I \to T\]
二番目の射は $T \otimes T$ のテンソル積を $T$ に写す。テンソル積は自己関手の合成によって与えられるので：
\[\mu \Colon T \circ T \to T\]
を得る。

\begin{figure}[H]
  \centering
  \includegraphics[width=0.3\textwidth]{images/monad.png}
\end{figure}

\noindent
これらはモナドを定義する二つの操作として認識できる（Haskell では \code{return} と \code{join} と呼ばれる）。モノイドの法則がモナドの法則に変わることも知っている。

では、この定義から自己関手についての言及をすべて取り除いてみよう。双圏 $\cat{C}$ から始めて、その中の $0$-セル $a$ を選ぶ。先に見たように、hom圏 $\cat{C}(a, a)$ はモノイダル圏だ。したがって、$1$-セル $T$ と二つの $2$-セル：
\begin{align*}
  \eta & \Colon I \to T         \\
  \mu  & \Colon T \circ T \to T
\end{align*}
を選ぶことで、$\cat{C}(a, a)$ におけるモノイドを定義できる。これがモノイドの法則を満たすとき、\emph{これを}モナドと呼ぶ。

\begin{figure}[H]
  \centering
  \includegraphics[width=0.3\textwidth]{images/bimonad.png}
\end{figure}

\noindent
これは $0$-セル、$1$-セル、$2$-セルのみを使ったモナドのはるかに一般的な定義だ。双圏 $\Cat$ に適用すると通常のモナドに帰着する。しかし他の双圏では何が起こるか見てみよう。

$\cat{Span}$ においてモナドを構成しよう。$0$-セルを選ぶ。それは集合で、すぐに理由が明らかになるが、$\mathit{Ob}$ と呼ぶことにする。次に自己 $1$-セルを選ぶ：$\mathit{Ob}$ から $\mathit{Ob}$ へのスパンだ。頂点に集合があり、$\mathit{Arr}$ と呼ぶ。二つの関数を持つ：
\begin{align*}
  \mathit{dom} & \Colon \mathit{Arr} \to \mathit{Ob} \\
  \mathit{cod} & \Colon \mathit{Arr} \to \mathit{Ob}
\end{align*}

\begin{figure}[H]
  \centering
  \includegraphics[width=0.3\textwidth]{images/spanmonad.png}
\end{figure}

\noindent
集合 $\mathit{Arr}$ の要素を「矢」と呼ぼう。$\mathit{Ob}$ の要素を「対象」と呼ぶよう言われれば、これがどこに向かっているかヒントが得られるかもしれない。二つの関数 $\mathit{dom}$ と $\mathit{cod}$ は「矢」に定義域と余定義域を割り当てる。

スパンをモナドにするためには、二つの $2$-セル $\eta$ と $\mu$ が必要だ。モノイダル単位は、この場合 $\mathit{Ob}$ を頂点とし二つの恒等関数を持つ $\mathit{Ob}$ から $\mathit{Ob}$ への自明なスパンだ。$2$-セル $\eta$ は頂点 $\mathit{Ob}$ と $\mathit{Arr}$ の間の関数だ。言い換えれば、$\eta$ はすべての「対象」に「矢」を割り当てる。$\cat{Span}$ における $2$-セルは可換条件を満たさなければならない——この場合：
\begin{align*}
  \mathit{dom} & \circ \eta = \id \\
  \mathit{cod} & \circ \eta = \id
\end{align*}

\begin{figure}[H]
  \centering
  \includegraphics[width=0.4\textwidth]{images/spanunit.png}
\end{figure}

\noindent
成分で書くと：
\[\mathit{dom}\ (\eta\ \mathit{ob}) = \mathit{ob} = \mathit{cod}\ (\eta\ \mathit{ob})\]
ここで $\mathit{ob}$ は $\mathit{Ob}$ の「対象」だ。言い換えれば、$\eta$ はすべての「対象」に、その定義域と余定義域がその「対象」である「矢」を割り当てる。この特別な「矢」を「恒等矢」と呼ぶことにする。

二番目の $2$-セル $\mu$ はスパン $\mathit{Arr}$ のそれ自身との合成に作用する。合成は引き戻しとして定義されるので、その要素は $\mathit{Arr}$ からの要素のペア——「矢」のペア $(a_1, a_2)$ だ。引き戻し条件は：
\[\mathit{cod}\ a_1 = \mathit{dom}\ a_2\]
一方の定義域がもう一方の余定義域であるとき、$a_1$ と $a_2$ は「合成可能」と言う。

\begin{figure}[H]
  \centering
  \includegraphics[width=0.5\textwidth]{images/spanmul.png}
\end{figure}

\noindent
$2$-セル $\mu$ は合成可能な「矢」のペア $(a_1, a_2)$ を $\mathit{Arr}$ の単一の「矢」$a_3$ に写す関数だ。言い換えれば $\mu$ は「矢」の合成を定義する。

モナドの法則が矢の単位律と結合律に対応することは容易に確認できる。私たちはちょうど圏（対象と矢が集合を形成する小圏）を定義したのだ。

要約すると、圏はスパンの双圏におけるモナドに他ならない。

この結果の驚くべき点は、圏をモナドやモノイドのような他の代数的構造と同じ土台に置くことだ。圏であることに特別なことは何もない。それはただ二つの集合と四つの関数だ。実際、対象の別個の集合さえ必要ない。なぜなら対象は恒等矢と同一視できるからだ（それらは一対一に対応している）。だからそれは本当にただの一つの集合といくつかの関数に過ぎない。数学全体において圏論が果たす中心的な役割を考えると、これは非常に謙虚な気づきだ。

\section{課題}

\begin{enumerate}
  \tightlist
  \item
        双圏における自己 $1$-セルの合成として定義されるテンソル積の単位律と結合律を導出せよ。
  \item
        $\cat{Span}$ におけるモナドのモナド法則が、結果として得られる圏の単位律と結合律に対応することを確認せよ。
  \item
        $\cat{Prof}$ におけるモナドが対象上の恒等関手であることを示せ。
  \item
        $\cat{Span}$ におけるモナドのモナド代数とは何か？
\end{enumerate}

\section{参考文献}
\begin{enumerate}
  \tightlist
  \item
        \urlref{https://graphicallinearalgebra.net/2017/04/16/a-monoid-is-a-category-a-category-is-a-monad-a-monad-is-a-monoid/}{Paweł Sobociński のブログ}。
\end{enumerate}
